\documentclass{article}
\usepackage[utf8]{inputenc}
\usepackage[T1]{fontenc}
\usepackage{geometry}
\geometry{
    a4paper,
    left=2.5cm,
    right=2.5cm,
    top=2.5cm,
    bottom=2.5cm
}
\usepackage{graphicx}
\usepackage{float}
\usepackage{hyperref}
\usepackage[english]{babel}
\usepackage{amsmath}
\usepackage{amssymb}
\usepackage{array}
\usepackage{booktabs}
\usepackage{caption}
\usepackage{subcaption}
\usepackage{enumitem}
\usepackage{fancyhdr}
\usepackage{titlesec}
\usepackage{url}
\usepackage{textcomp}
\usepackage{xcolor}
\usepackage{listings}
\usepackage{multirow}
\usepackage{longtable}
\usepackage{tabularx}
\usepackage{times}

% Set up page style
\pagestyle{fancy}
\fancyhf{}
\rfoot{\thepage}
\renewcommand{\headrulewidth}{0pt}

% Title formatting
\titleformat{\chapter}[display]{\normalfont\bfseries}{\chaptertitle\thechapter}{0pt}{\Huge}
\titlespacing*{\chapter}{0pt}{50pt}{40pt}

% Document info
\title{Template Document}
\author{Academic Support Agent}
\date{\today}

\begin{document}
\maketitle

% Seleciona o idioma do documento (conforme pacotes do babel)
%\selectlanguage{english}
\selectlanguage{brazil}

% Retira espa o extra obsoleto entre as frases.
\frenchspacing 

% ----------------------------------------------------------
% ELEMENTOS PR -TEXTUAIS
% ----------------------------------------------------------
% \pretextual

% ---
% Capa
% ---
\imprimircapa
% ---

% ---
% Folha de rosto
% (o * indica que haver  a ficha bibliogr fica)
% ---
\imprimirfolhaderosto*
% ---

% ---
% Inserir a ficha bibliografica
% ---

% Isto   um exemplo de Ficha Catalogr fica, ou ``Dados internacionais de
% cataloga o-na-publica o''. Voc  pode utilizar este modelo como refer ncia. 
% Por m, provavelmente a biblioteca da sua universidade lhe fornecer  um PDF
% com a ficha catalogr fica definitiva ap s a defesa do trabalho. Quando estiver
% com o documento, salve-o como PDF no diret rio do seu projeto e substitua todo
% o conte do de implementa o deste arquivo pelo comando abaixo:
%
% \begin{fichacatalografica}
%     \includepdf{fig_ficha_catalografica.pdf}
% \end{fichacatalografica}

\begin{fichacatalografica}
	\sffamily
	\vspace*{\fill}					% Posi o vertical
	\begin{center}					% Minipage Centralizado
	\fbox{\begin{minipage}[c][8cm]{13.5cm}		% Largura
	\small
	Solicite sua ficha catalogr fica em: \url{http://bcficat.ufpa.br/}
	\end{minipage}}
	\end{center}
\end{fichacatalografica}
% ---

% ---
% Inserir errata
% ---
%\begin{errata}
%Elemento opcional da %\cite{NBR14724:2011}. Exemplo:

%\vspace{\onelineskip}

%FERRIGNO, C. R. A. \textbf{Tratamento de neoplasias  sseas apendiculares com
%reimplanta o de enxerto  sseo aut logo autoclavado associado ao plasma
%rico em plaquetas}: estudo cr tico na cirurgia de preserva o de membro em
%c es. 2011. 128 f. Tese (Livre-Doc ncia) - Faculdade de Medicina Veterin ria e
%Zootecnia, Universidade de S o Paulo, S o Paulo, 2011.

%\begin{table}[htb]
%\center
%\footnotesize
%\begin{tabular}{|p{1.4cm}|p{1cm}|p{3cm}|p{3cm}|}
%  \hline
%   \textbf{Folha} & \textbf{Linha}  & \textbf{Onde se l }  & \textbf{Leia-se}  \\
%    \hline
%    1 & 10 & auto-conclavo & autoconclavo\\
%   \hline
%\end{tabular}
%\end{table}

%\end{errata}
% ---

% ---
% Inserir folha de aprova o
% ---

% Isto   um exemplo de Folha de aprova o, elemento obrigat rio da NBR
% 14724/2011 (se o 4.2.1.3). Voc  pode utilizar este modelo at  a aprova o
% do trabalho. Ap s isso, substitua todo o conte do deste arquivo por uma
% imagem da p gina assinada pela banca com o comando abaixo:
%
% \begin{folhadeaprovacao}
% \includepdf{folhadeaprovacao_final.pdf}
% \end{folhadeaprovacao}
%
\begin{folhadeaprovacao}

  \begin{center}
    {\ABNTEXchapterfont\large\imprimirautor}

    \vspace*{\fill}\vspace*{\fill}
    \begin{center}
      \ABNTEXchapterfont\bfseries\Large\imprimirtitulo
    \end{center}
    \vspace*{\fill}
    
    \hspace{.45\textwidth}
    \begin{minipage}{.5\textwidth}
        \imprimirpreambulo
    \end{minipage}%
    \vspace*{\fill}
   \end{center}
        
   Conceito: \rule{3cm}{.1pt}
   
   \imprimirlocal, 1 de janeiro de 2019.
   
   \vspace{1cm}
   \begin{center}
   BANCA EXAMINADORA
   \end{center}
    

   \assinatura{\textbf{\imprimirorientador} - Orientador \\ UFPA}
   %\assinatura{\textbf{\imprimircoorientador} - Coorientador \\ UFPA}
   \assinatura{\textbf{Nome Convidado 1} \\ SIGLA INSTITUI O}
   \assinatura{\textbf{Nome Convidado 2} \\ SIGLA INSTITUI O}
   %\assinatura{\textbf{Nome Convidado 3} \\ SIGLA INSTITUI O}
      

  
\end{folhadeaprovacao}
% ---

% ---
% Dedicat ria
% ---
\begin{dedicatoria}
   \vspace*{\fill}
   \centering
   \noindent
   \textit{ Escreva sua dedicat ria aqui.} \vspace*{\fill}
\end{dedicatoria}
% ---

% ---
% Agradecimentos
% ---
\begin{agradecimentos}
Os agradecimentos principais s o direcionados   Gerald Weber, Miguel Frasson,
Leslie H. Watter, Bruno Parente Lima, Fl vio de Vasconcellos Corr a, Otavio Real
Salvador, Renato Machnievscz e todos aqueles que
contribu ram para que a produ o de trabalhos acad micos conforme
as normas ABNT com \LaTeX\ fosse poss vel.

Agradecimentos especiais s o direcionados ao Centro de Pesquisa em Arquitetura
da Informa o da Universidade de
Bras lia (CPAI), ao grupo de usu rios
\emph{latex-br} e aos
novos volunt rios do grupo
\emph{\abnTeX} e que contribu ram e que ainda
contribuir o para a evolu o do \abnTeX.

\end{agradecimentos}
% ---

% ---
% Ep grafe
% ---
\begin{epigrafe}
    \vspace*{\fill}
	\begin{flushright}
		\textit{``Escreva sua ep grafe aqui''\\
		(Fulano de Tal, 19XX)}
	\end{flushright}
\end{epigrafe}
% ---

% ---
% RESUMOS
% ---

% resumo em portugu s
\setlength{\absparsep}{18pt} % ajusta o espa amento dos par grafos do resumo
\begin{resumo}
 Segundo a \cite{NBR6028:2003}, o resumo deve ressaltar o
 objetivo, o m todo, os resultados e as conclus es do documento. A ordem e a extens o
 destes itens dependem do tipo de resumo (informativo ou indicativo) e do
 tratamento que cada item recebe no documento original. O resumo deve ser
 precedido da refer ncia do documento, com exce o do resumo inserido no
 pr prio documento. (\ldots) As palavras-chave devem figurar logo abaixo do
 resumo, antecedidas da express o Palavras-chave:, separadas entre si por
 ponto e finalizadas tamb m por ponto.

 \textbf{Palavras-chave}: latex. abntex. editora o de texto.
\end{resumo}

% resumo em ingl s
\begin{resumo}[Abstract]
 \begin{otherlanguage*}{english}
   This is the english abstract.

   \vspace{\onelineskip}
 
   \noindent 
   \textbf{Keywords}: latex. abntex. text editoration.
 \end{otherlanguage*}
\end{resumo}

% ---

% ---
% inserir lista de ilustra es
% ---
\pdfbookmark[0]{\listfigurename}{lof}


% ---

% ---
% inserir lista de quadros
% ---
\pdfbookmark[0]{\listofquadrosname}{loq}
\listofquadros*

% ---

% ---
% inserir lista de tabelas
% ---
\pdfbookmark[0]{\listtablename}{lot}


% ---

% ---
% inserir lista de abreviaturas e siglas
% ---
\begin{siglas}
  \item[ABNT] Associa o Brasileira de Normas T cnicas
  \item[abnTeX] ABsurdas Normas para TeX
\end{siglas}
% ---

% ---
% inserir lista de s mbolos
% ---
\begin{simbolos}
  \item[$ \Gamma $] Letra grega Gama
  \item[$ \Lambda $] Lambda
  \item[$ \zeta $] Letra grega min scula zeta
  \item[$ \in $] Pertence
\end{simbolos}
% ---

% ---
% inserir o sumario
% ---
\pdfbookmark[0]{\contentsname}{toc}


% ---



% ----------------------------------------------------------
% ELEMENTOS TEXTUAIS
% ----------------------------------------------------------
\textual

% ----------------------------------------------------------
% Introdu o (exemplo de cap tulo sem numera o, mas presente no Sum rio)
% ----------------------------------------------------------
\chapter{Introdu o}
% ----------------------------------------------------------

Este documento e seu c digo-fonte s o exemplos de refer ncia de uso da classe
\textsf{abntex2} e do pacote \textsf{abntex2cite}. O documento 
exemplifica a elabora o de trabalho acad mico (tese, disserta o e outros do
g nero) produzido conforme a ABNT NBR 14724:2011 \emph{Informa o e documenta o
- Trabalhos acad micos - Apresenta o}.

A express o ``Modelo Can nico''   utilizada para indicar que \abnTeX\ n o  
modelo espec fico de nenhuma universidade ou institui o, mas que implementa t o
somente os requisitos das normas da ABNT. Uma lista completa das normas
observadas pelo \abnTeX\   apresentada em \cite{abntex2classe}.

Sinta-se convidado a participar do projeto \abnTeX! Acesse o site do projeto em
\url{http://www.abntex.net.br/}. Tamb m fique livre para conhecer,
estudar, alterar e redistribuir o trabalho do \abnTeX, desde que os arquivos
modificados tenham seus nomes alterados e que os cr ditos sejam dados aos
autores originais, nos termos da ``The \LaTeX\ Project Public
License''\footnote{\url{http://www.latex-project.org/lppl.txt}}.

Encorajamos que sejam realizadas customiza es espec ficas deste exemplo para
universidades e outras institui es --- como capas, folha de aprova o, etc.
Por m, recomendamos que ao inv s de se alterar diretamente os arquivos do
\abnTeX, distribua-se arquivos com as respectivas customiza es.
Isso permite que futuras vers es do \abnTeX~n o se tornem automaticamente
incompat veis com as customiza es promovidas. Consulte
\cite{abntex2-wiki-como-customizar} para mais informa es.

Este documento deve ser utilizado como complemento dos manuais do \abnTeX\ 
\cite{abntex2classe,abntex2cite,abntex2cite-alf} e da classe \textsf{memoir}
\cite{memoir}. 

Esperamos, sinceramente, que o \abnTeX\ aprimore a qualidade do trabalho que
voc  produzir , de modo que o principal esfor o seja concentrado no principal:
na contribui o cient fica.

Equipe \abnTeX 

Lauro C sar Araujo

% ---
% Capitulo com exemplos de comandos inseridos de arquivo externo 
% ---
% Content from abntex2-modelo-include-comandos would be included here

% ---

\chapter{Conte dos espec ficos do modelo de trabalho acad mico}\label{cap_trabalho_academico}

\section{Quadros}

Este modelo vem com o ambiente \texttt{quadro} e impress o de Lista de quadros 
configurados por padr o. Verifique um exemplo de utiliza o:

\begin{quadro}[htb]
\caption{\label{quadro_exemplo}Exemplo de quadro}
\begin{tabular}{|c|c|c|c|}
	\hline
	\textbf{Pessoa} & \textbf{Idade} & \textbf{Peso} & \textbf{Altura} \\ \hline
	Marcos & 26    & 68   & 178    \\ \hline
	Ivone  & 22    & 57   & 162    \\ \hline
	...    & ...   & ...  & ...    \\ \hline
	Sueli  & 40    & 65   & 153    \\ \hline
\end{tabular}
\fonte{Autor.}
\end{quadro}

Este par grafo apresenta como referenciar o quadro no texto, requisito
obrigat rio da ABNT. 
Primeira op o, utilizando \texttt{autoref}: Ver o \autoref{quadro_exemplo}. 
Segunda op o, utilizando  \texttt{ref}: Ver o Quadro \ref{quadro_exemplo}.


% ---
% Capitulo de revis o de literatura
% ---
\chapter{Lorem ipsum dolor sit amet}
% ---

% ---
\section{Aliquam vestibulum fringilla lorem}
% ---

\lipsum[1]

\lipsum[2-3]

% ---
% primeiro capitulo de Resultados
% ---
\chapter{Lectus lobortis condimentum}
% ---

% ---
\section{Vestibulum ante ipsum primis in faucibus orci luctus et ultrices
posuere cubilia Curae}
% ---

\lipsum[21-22]

% ---
% segundo capitulo de Resultados
% ---
\chapter{Nam sed tellus sit amet lectus urna ullamcorper tristique interdum
elementum}
% ---

% ---
\section{Pellentesque sit amet pede ac sem eleifend consectetuer}
% ---

\lipsum[24]

% ----------------------------------------------------------
% Finaliza a parte no bookmark do PDF
% para que se inicie o bookmark na raiz
% e adiciona espa o de parte no Sum rio
% ----------------------------------------------------------
\phantompart

% ---
% Conclus o
% ---
\chapter{Conclus o}
% ---

\lipsum[31-33]

% ----------------------------------------------------------
% ELEMENTOS P S-TEXTUAIS
% ----------------------------------------------------------
\postextual
% ----------------------------------------------------------

% ----------------------------------------------------------
% Refer ncias bibliogr ficas
% ----------------------------------------------------------


% ----------------------------------------------------------
% Gloss rio
% ----------------------------------------------------------
%
% Consulte o manual da classe abntex2 para orienta es sobre o gloss rio.
%
%\glossary

% ----------------------------------------------------------
% Ap ndices
% ----------------------------------------------------------

% ---
% Inicia os ap ndices
% ---
\begin{apendicesenv}

% Imprime uma p gina indicando o in cio dos ap ndices
%\partapendices

% ----------------------------------------------------------
\chapter{Quisque libero justo}
% ----------------------------------------------------------

\lipsum[50]

% ----------------------------------------------------------
\chapter{Nullam elementum urna vel imperdiet sodales elit ipsum pharetra ligula
ac pretium ante justo a nulla curabitur tristique arcu eu metus}
% ----------------------------------------------------------
\lipsum[55-57]

\end{apendicesenv}
% ---


% ----------------------------------------------------------
% Anexos
% ----------------------------------------------------------

% ---
% Inicia os anexos
% ---
\begin{anexosenv}

% Imprime uma p gina indicando o in cio dos anexos
%\partanexos

% ---
\chapter{Morbi ultrices rutrum lorem.}
% ---
\lipsum[30]

% ---
\chapter{Cras non urna sed feugiat cum sociis natoque penatibus et magnis dis
parturient montes nascetur ridiculus mus}
% ---

\lipsum[31]

% ---
\chapter{Fusce facilisis lacinia dui}
% ---

\lipsum[32]

\end{anexosenv}
\end{document}